\documentclass{beamer}
\usetheme{metropolis} % or 'madrid', 'default'—adds clean structure

\title{Definitions, Theorems, Proofs, and Proof Writing}
\author{Math Club}
\date{\today}

\begin{document}

\frame{\titlepage}

\begin{frame}{Motivation}
  \textbf{Question:} Is the sum of two odd numbers even or odd?

  \vspace{0.5em}
  \textbf{Examples:}
  \[
    2 + 3 = 4 \qquad 7 + 5 = 13 \qquad 5 + 1 = 6 \qquad -3 + 1 = -2
  \]
  
  \vspace{0.5em}
  With further examples we see that the pattern holds
  
  \vspace{1em}
  How can we prove this is true for \emph{all} integers without testing each case?
\end{frame}

\begin{frame}{Theorems}
    First, we must explicitly write what we are trying to prove.
    A \textbf{theorem} is a mathematical statement that has been proven to be true on the basis of previously established statements.
    \vspace{0.5em}
    \textbf{Theorem:} The sum of two odd numbers is even.
    \vspace{0.5em}
    However, for rigor and completeness we must know exactly what an even or odd number is. 
\end{frame}

\begin{frame}{Definitions}
  A theorem must be based in precise and agreed upoon definitions.
  A \textbf{definition} establishes the exact, non-negotiable meaning of a term. It is either satifised or not.

  \vspace{1em}
    Definitions are \textbf{agreed upon} as a fact rather than proven. Our theorem and proofs are then built upon these.
\end{frame}

\begin{frame}{Definitions: Parity of Integers}

  \vspace{0.5em}
  Recall the set of integers is $\mathbb{Z} = \{ \dots, -2, -1, 0, 1, 2, \dots \}$.

  \vspace{0.5em}
  \begin{block}{Definition: Even and Odd Integers}
    Let $n \in \mathbb{Z}$.
    \begin{itemize}
      \item $n$ is \textbf{even} if $n = 2k$ for some integer $k \in \mathbb{Z}$.
      \item $n$ is \textbf{odd} if $n = 2k + 1$ for some integer $k \in \mathbb{Z}$.
    \end{itemize}
  \end{block}

  \vspace{0.5em}
  This definition:
  \begin{itemize}
    \item Sets the boundary to apply only to integers
    \item Paritions all integers into being either even or odd with no in between
    \item Provides us with variables and mathematical expressions that can be manipulated
  \end{itemize}
  \vspace{0.5em}
\end{frame}

\begin{frame}{The Logic of a Direct Proof}
  Most mathematical claims can be structured as conditional statements:
  \[
    P \implies Q \quad \text{``If P, then Q''}
  \]

  \begin{itemize}
    \item \textbf{Hypothesis ($P$):} What you are given as true (e.g., \emph{``$n$ and $m$ are odd integers''}).
    \item \textbf{Conclusion ($Q$):} What you want to deduce (e.g., \emph{``$n + m$ is an even integer''}).
  \end{itemize}

  \vspace{0.5em}
  \textbf{The Logical Chain:}
  A direct proof bridges $P$ to $Q$ through a sequence of valid deductions:
  \[P \implies A_1 \implies A_2 \implies \dots \implies A_n \implies Q\]

  $P$ is true by assumption, so if our logic is correct in each step, all following statements are true, thus leaving us with $Q$
\end{frame}

\begin{frame}{Scratchwork}
    Before writing a formal proof, test your ideas and map out the logic on scratch paper.
\vspace{0.5em}
  \textbf{Workflow}
  \begin{itemize}
    \item \textbf{Unpack Definitions:} Translate the assumptions ($P$) using definitions
    \item \textbf{Identify Target State:} Detirmine what the conclusion ($Q$) needs to look like structurally.
    \item \textbf{Manipulate \& Connect:} Apply algebra, logic, etc to reach $Q$.
  \end{itemize}
\end{frame}

\end{document}